% Source: https://artofproblemsolving.com/community/c3h1946190_hard_absolute_value_inequalities
% Stable identity: aops-topic-1946190-q01
The expressions inside the absolute values change sign at $x=-5$ and $x=\dfrac{7}{3}$, so consider three cases.
\begin{itemize}
\item $x<-5$: the inequality is $-(x+5)-(7-3x)>1$, i.e. $2x-12>1$, $x>\dfrac{13}{2}$. No solutions here.
\item $-5\le x<\dfrac{7}{3}$: it is $(x+5)-(7-3x)>1$, i.e. $4x-2>1$, $x>\dfrac{3}{4}$. Solutions $\dfrac{3}{4}<x<\dfrac{7}{3}$.
\item $x\ge\dfrac{7}{3}$: it is $(x+5)-(3x-7)>1$, i.e. $-2x+12>1$, $x<\dfrac{11}{2}$. Solutions $\dfrac{7}{3}\le x<\dfrac{11}{2}$.
\end{itemize}
Combining, the solution set is $\boxed{\dfrac{3}{4}<x<\dfrac{11}{2}}$. (Check: at $x=\dfrac{3}{4}$ and at $x=\dfrac{11}{2}$ the left side equals $1$.)
